\section{Inline Visual Route: Model Foundation}

These figures are placed in the main argument so the reader can learn the model before opening evidence tables.


\subsection{Carrier to realization}

This visual anchor names the objects, review direction, and formula or numeric boundary that keep the figure tied to the argument.


\begin{figure}[htbp]

\centering

\begin{tikzpicture}[>=Latex,node distance=2.4cm]

\node[draw,rounded corners,fill=blue!7,text width=0.28\textwidth,align=center,minimum height=1.0cm] (a) {Carrier \$C\$};

\node[draw,rounded corners,fill=green!7,text width=0.28\textwidth,align=center,minimum height=1.0cm,right=of a] (b) {Realization \$R\_t\$};

\draw[->,line width=0.8pt] (a) -- node[above,align=center] {lawful realization} (b);

\node[draw,rounded corners,fill=orange!8,text width=0.68\textwidth,align=center,below=0.9cm of $(a)!0.5!(b)$] (c) {review question: support class, boundary, falsifier};

\draw[->,dashed] (c.north) -- ($(a)!0.5!(b)$);

\end{tikzpicture}

\caption{Carrier to realization. Shows the basic OC route from carrier to time-indexed realization; numeric anchor: \$t=0,1\$ status comparison. The figure is placed inline in the manuscript; complete machine evidence remains in the evidence package.}

\label{fig:r005-inline-28a_oc133_inline_figures_foundation-1}

\end{figure}

\subsection{Liveness status split}

This visual anchor names the objects, review direction, and formula or numeric boundary that keep the figure tied to the argument.


\begin{figure}[htbp]

\centering

\begin{tikzpicture}[>=Latex,node distance=2.4cm]

\node[draw,rounded corners,fill=blue!7,text width=0.28\textwidth,align=center,minimum height=1.0cm] (a) {Live};

\node[draw,rounded corners,fill=green!7,text width=0.28\textwidth,align=center,minimum height=1.0cm,right=of a] (b) {Dead};

\draw[->,line width=0.8pt] (a) -- node[above,align=center] {status predicate} (b);

\node[draw,rounded corners,fill=orange!8,text width=0.68\textwidth,align=center,below=0.9cm of $(a)!0.5!(b)$] (c) {review question: support class, boundary, falsifier};

\draw[->,dashed] (c.north) -- ($(a)!0.5!(b)$);

\end{tikzpicture}

\caption{Liveness status split. Separates live status from dead status; formula anchor: \$live(x,t)\textbackslash{}in\textbackslash{}\{0,1\textbackslash{}\}\$. The figure is placed inline in the manuscript; complete machine evidence remains in the evidence package.}

\label{fig:r005-inline-28a_oc133_inline_figures_foundation-2}

\end{figure}

\subsection{Residue after collapse}

This visual anchor names the objects, review direction, and formula or numeric boundary that keep the figure tied to the argument.


\begin{figure}[htbp]

\centering

\begin{tikzpicture}[>=Latex,node distance=2.4cm]

\node[draw,rounded corners,fill=blue!7,text width=0.28\textwidth,align=center,minimum height=1.0cm] (a) {Collapse};

\node[draw,rounded corners,fill=green!7,text width=0.28\textwidth,align=center,minimum height=1.0cm,right=of a] (b) {Residue};

\draw[->,line width=0.8pt] (a) -- node[above,align=center] {trace} (b);

\node[draw,rounded corners,fill=orange!8,text width=0.68\textwidth,align=center,below=0.9cm of $(a)!0.5!(b)$] (c) {review question: support class, boundary, falsifier};

\draw[->,dashed] (c.north) -- ($(a)!0.5!(b)$);

\end{tikzpicture}

\caption{Residue after collapse. Shows why residue is not continuing liveness; numeric anchor: continuumness score reaches \$0\$. The figure is placed inline in the manuscript; complete machine evidence remains in the evidence package.}

\label{fig:r005-inline-28a_oc133_inline_figures_foundation-3}

\end{figure}

\subsection{Boundary crossing}

This visual anchor names the objects, review direction, and formula or numeric boundary that keep the figure tied to the argument.


\begin{figure}[htbp]

\centering

\begin{tikzpicture}[>=Latex,node distance=2.4cm]

\node[draw,rounded corners,fill=blue!7,text width=0.28\textwidth,align=center,minimum height=1.0cm] (a) {Interior};

\node[draw,rounded corners,fill=green!7,text width=0.28\textwidth,align=center,minimum height=1.0cm,right=of a] (b) {Boundary};

\draw[->,line width=0.8pt] (a) -- node[above,align=center] {\$\textbackslash{}partial\textbackslash{}Omega\$} (b);

\node[draw,rounded corners,fill=orange!8,text width=0.68\textwidth,align=center,below=0.9cm of $(a)!0.5!(b)$] (c) {review question: support class, boundary, falsifier};

\draw[->,dashed] (c.north) -- ($(a)!0.5!(b)$);

\end{tikzpicture}

\caption{Boundary crossing. Shows a threshold boundary as a reviewable object rather than a metaphor; formula anchor: \$d(x,\textbackslash{}partial\textbackslash{}Omega)\$. The figure is placed inline in the manuscript; complete machine evidence remains in the evidence package.}

\label{fig:r005-inline-28a_oc133_inline_figures_foundation-4}

\end{figure}

\subsection{Operator stack}

This visual anchor names the objects, review direction, and formula or numeric boundary that keep the figure tied to the argument.


\begin{figure}[htbp]

\centering

\begin{tikzpicture}[>=Latex,node distance=2.4cm]

\node[draw,rounded corners,fill=blue!7,text width=0.28\textwidth,align=center,minimum height=1.0cm] (a) {Flow};

\node[draw,rounded corners,fill=green!7,text width=0.28\textwidth,align=center,minimum height=1.0cm,right=of a] (b) {Guard reset};

\draw[->,line width=0.8pt] (a) -- node[above,align=center] {hybrid update} (b);

\node[draw,rounded corners,fill=orange!8,text width=0.68\textwidth,align=center,below=0.9cm of $(a)!0.5!(b)$] (c) {review question: support class, boundary, falsifier};

\draw[->,dashed] (c.north) -- ($(a)!0.5!(b)$);

\end{tikzpicture}

\caption{Operator stack. Separates smooth flow from discrete update; formula anchor: \$x\_\{t+1\}=G(F\_t(x\_t))\$. The figure is placed inline in the manuscript; complete machine evidence remains in the evidence package.}

\label{fig:r005-inline-28a_oc133_inline_figures_foundation-5}

\end{figure}

\subsection{K-level witness}

This visual anchor names the objects, review direction, and formula or numeric boundary that keep the figure tied to the argument.


\begin{figure}[htbp]

\centering

\begin{tikzpicture}[>=Latex,node distance=2.4cm]

\node[draw,rounded corners,fill=blue!7,text width=0.28\textwidth,align=center,minimum height=1.0cm] (a) {Lower projection};

\node[draw,rounded corners,fill=green!7,text width=0.28\textwidth,align=center,minimum height=1.0cm,right=of a] (b) {Retained witness};

\draw[->,line width=0.8pt] (a) -- node[above,align=center] {non-loss test} (b);

\node[draw,rounded corners,fill=orange!8,text width=0.68\textwidth,align=center,below=0.9cm of $(a)!0.5!(b)$] (c) {review question: support class, boundary, falsifier};

\draw[->,dashed] (c.north) -- ($(a)!0.5!(b)$);

\end{tikzpicture}

\caption{K-level witness. Shows the reduction test for a K-level claim; numeric anchor: witness preserved or lost as \$0/1\$. The figure is placed inline in the manuscript; complete machine evidence remains in the evidence package.}

\label{fig:r005-inline-28a_oc133_inline_figures_foundation-6}

\end{figure}
